%%%%%%%%%%%%%%%%%%%%%%%%%%%%%%%%%%%%%%%%%%%%%%%%%%%%%%%%%%%%%%
%% Web GIS Laboratory: District Explorer
%% Institute of Geographical Information Systems, NUST
%%
%% Layout follows "A Psychology Lab Manual" by John Paul Minda
%% (Overleaf, CC BY 4.0): article class, fullpage margins,
%% Helvetica throughout, parskip, blue links, no colour boxes.
%%
%% COMPILER: pdfLaTeX
%%%%%%%%%%%%%%%%%%%%%%%%%%%%%%%%%%%%%%%%%%%%%%%%%%%%%%%%%%%%%%

\documentclass{article}
\usepackage{fullpage}

\renewcommand{\familydefault}{\sfdefault}
\usepackage[scaled=1]{helvet}
\usepackage[helvet]{sfmath}
\everymath={\sf}

\usepackage{parskip}
\usepackage[colorlinks=true, allcolors=blue]{hyperref}

%% Required by this document but not by the original template.
%% Encoding first: without T1, < and > inside code print as
%% inverted punctuation and straight quotes turn curly.
\usepackage{iftex}
\ifPDFTeX
  \usepackage[T1]{fontenc}
  \usepackage[utf8]{inputenc}
\fi
\usepackage{textcomp}
\usepackage{upquote}
\usepackage{xcolor}
\usepackage{listings}
\usepackage{booktabs}
\usepackage{tabularx}
\usepackage{array}
\usepackage{enumitem}

\setlist[enumerate]{itemsep=.25em,topsep=.5em}
\setlist[itemize]{itemsep=.25em,topsep=.5em}
\renewcommand{\arraystretch}{1.25}
\newcolumntype{L}[1]{>{\raggedright\arraybackslash}p{#1}}

%% Inline code. \cd for short tokens, \cdb for long paths and URLs,
%% which must be allowed to break across lines.
\urlstyle{tt}
\newcommand{\cd}[1]{\texttt{#1}}
\newcommand{\cdb}[1]{\nolinkurl{#1}}

%% Code blocks: a plain rule, no shading, in keeping with the template
\lstdefinestyle{plain}{
  basicstyle=\ttfamily\small,
  frame=single,
  rulecolor=\color{black!35},
  framesep=7pt,
  xleftmargin=7pt,
  xrightmargin=7pt,
  breaklines=true,
  columns=fullflexible,
  keepspaces=true,
  showstringspaces=false,
  upquote=true,
  aboveskip=1em,
  belowskip=1em,
}
\lstset{style=plain}
\lstnewenvironment{code}{\lstset{style=plain}}{}

%% Caption for a code block. Deliberately not \paragraph, which is a
%% run-in heading and would be overlapped by the listing frame.
\newcommand{\codecap}[1]{%
  \par\vspace{.9em}\noindent\textbf{#1}\par\vspace{-.6em}}

%% A block that must not be skimmed past: ruled above and below,
%% with a bold label. Used for warnings and checkpoints.
\newenvironment{keypoint}[1]
  {\par\vspace{.7em}\noindent\rule{\linewidth}{.6pt}\par\vspace{.15em}
   \noindent\textbf{#1}\par\nobreak\vspace{-.3em}}
  {\par\vspace{.15em}\noindent\rule{\linewidth}{.6pt}\par\vspace{.7em}}

\title{Web GIS Laboratory\\[.3em]\large District Explorer:
       building a web map from a shapefile}
\author{Dr.\ Shahid Nawaz Khan\\
        Institute of Geographical Information Systems\\
        School of Civil and Environmental Engineering\\
        National University of Sciences and Technology, Islamabad}
\date{MS Remote Sensing and GIS}

\setcounter{tocdepth}{2}

\begin{document}
\maketitle

\begin{tabularx}{\textwidth}{@{}L{3.4cm} X@{}}
\toprule
Duration    & 3 hours \\
Covers      & Both preceding lectures \\
Software    & QGIS, Visual Studio Code, a browser, Git \\
Deliverable & A published map and a measurement report \\
\bottomrule
\end{tabularx}

\tableofcontents


%==============================================================
\section{Before you begin}

\subsection{What you will build}

A working web map of Pakistan's districts, published at a public URL.
The finished page lets a user pick a province from a drop-down, redraws
the map with only those districts, shades each district by population,
and shows its attributes when clicked.

You will also produce a short measurement report comparing file sizes,
load times and computed areas, using your own numbers rather than the
figures quoted in the lecture.

\subsection{Learning outcomes}

By the end of this laboratory you will be able to:

\begin{enumerate}
  \item Convert a shapefile to GeoJSON and control its size.
  \item Explain the effect of coordinate precision and simplification on
        file size, and choose values appropriate to a map's scale.
  \item Read a GeoJSON file as text and identify its structure and its
        errors.
  \item Request a file from JavaScript and inspect it in the browser
        console.
  \item Draw, style, classify and filter a vector layer with Leaflet.
  \item Connect an HTML control to a map through an event listener.
  \item Measure the difference between an area computed in Web Mercator
        and one computed in a projected system.
  \item Publish a static site and read its network activity.
\end{enumerate}

\subsection{Structure and timing}

Parts 1 to 7 are the core of the laboratory and everyone is expected to
finish them. Parts 8 and 9 are extensions. If you run short of time,
complete Part 10 first so that your work is published, then return to
the extensions.

\vspace{.6em}
\begin{tabularx}{\textwidth}{@{}c X r@{}}
\toprule
\textbf{Part} & \textbf{What you do} & \textbf{Time} \\
\midrule
1  & Set up the project                          & 15 min \\
2  & Prepare the data in QGIS                    & 30 min \\
3  & Read the data as text                       & 15 min \\
4  & Build the page                              & 20 min \\
5  & Load the data with JavaScript               & 20 min \\
6  & Draw it with Leaflet                        & 25 min \\
7  & Classify and add a popup                    & 20 min \\
\midrule
8  & Filter with a control \textit{(extension)}  & 25 min \\
9  & Measure the coordinate systems \textit{(extension)} & 20 min \\
\midrule
10 & Publish and submit                          & 20 min \\
\bottomrule
\end{tabularx}

\subsection{What you need}

\begin{tabularx}{\textwidth}{@{}L{4.2cm} X@{}}
\toprule
\textbf{Item} & \textbf{Requirement} \\
\midrule
QGIS            & Version 3.28 or later \\
Visual Studio Code & With the Live Server extension installed \\
Browser         & Chrome or Edge, for the developer tools \\
Git             & Installed and configured with your name and email \\
Account         & A GitHub account \\
Data            & The districts shapefile from LMS, about 9 MB zipped \\
\bottomrule
\end{tabularx}

\begin{keypoint}{Field names in your data}
This handout uses the field names \cd{district}, \cd{province},
\cd{pop\_2023} and \cd{area\_km2}. Open the attribute table in QGIS
before you start and check what your copy actually uses. If the names
differ, substitute yours everywhere they appear in the code.
\end{keypoint}


%==============================================================
\section{Set up the project}

\subsection{Create the folder}

Open PowerShell and run:

\begin{code}
cd C:\
mkdir webgis-lab
cd C:\webgis-lab
mkdir data
code .
\end{code}

\begin{keypoint}{Where the folder goes}
Use \cd{C:\textbackslash webgis-lab}. Do not use Documents, Desktop, or
any folder inside OneDrive. Those folders are synced, and the sync
process locks files while a web server is reading them. Spaces in the
path cause a separate set of problems later.
\end{keypoint}


\subsection{Create the repository}

\begin{enumerate}
  \item On \cd{github.com}, create a new repository named
        \cd{webgis-lab-<your-roll-number>}.
  \item Set it to \textbf{Public}. This is required for free hosting.
  \item Do not tick any of the initialise options.
  \item Connect your folder to it:
\end{enumerate}

\begin{code}
git init
git branch -M main
git remote add origin https://github.com/<your-username>/webgis-lab-<roll>.git
\end{code}

\subsection{Create the three files}

In the VS Code explorer, create these in \cd{C:\textbackslash webgis-lab}:

\begin{tabularx}{\textwidth}{@{}L{3.4cm} X@{}}
\toprule
\textbf{File} & \textbf{Holds} \\
\midrule
\cd{index.html} & The structure of the page \\
\cd{style.css}  & How the page looks \\
\cd{app.js}     & What the page does \\
\bottomrule
\end{tabularx}

\vspace{.8em}
Keeping the three languages in three files is the normal arrangement
once a page grows beyond a few lines. The HTML file will link to the
other two.

\begin{keypoint}{Checkpoint 1}
A screenshot of VS Code showing the folder \cd{C:\textbackslash
webgis-lab} with \cd{index.html}, \cd{style.css}, \cd{app.js} and an
empty \cd{data} folder.
\end{keypoint}


%==============================================================
\section{Prepare the data in QGIS}

This part produces four versions of the same layer. Record the size of
each one as you go, because the comparison is part of your report.

\subsection{Inspect what you were given}

\begin{enumerate}
  \item Unzip the shapefile into a working folder \textbf{outside}
        \cd{C:\textbackslash webgis-lab}. Raw data is not committed to
        the repository.
  \item Load it into QGIS and open the attribute table. Note the field
        names and their types.
  \item Open \textbf{Layer Properties}, then \textbf{Information}, and
        record the coordinate reference system and the feature count.
  \item In your file manager, note the total size of all the shapefile
        components together.
\end{enumerate}

\paragraph{Why the field types matter}
A population column stored as text rather than as a number will sort and
compare alphabetically, so 9 will appear greater than 1500. This breaks
any classification built on it and produces no error. Check the type now
rather than after your choropleth looks wrong.


\subsection{Version A: a straight conversion}

Right-click the layer, choose \textbf{Export}, then \textbf{Save
Features As}.

\begin{tabularx}{\textwidth}{@{}L{4.6cm} X@{}}
\toprule
\textbf{Setting} & \textbf{Value} \\
\midrule
Format              & GeoJSON \\
File name           & \cd{districts\_A.geojson} \\
CRS                 & EPSG:4326, WGS 84 \\
Coordinate precision & Leave at the default \\
\bottomrule
\end{tabularx}

\vspace{.8em}
Record the resulting file size.

\begin{keypoint}{Always export in EPSG:4326}
GeoJSON assumes longitude and latitude in WGS 84. Exporting in any other
system produces a file that is syntactically valid and silently placed in
the wrong part of the world. If your data arrived in a projected system,
this export step is also the reprojection step.
\end{keypoint}


\subsection{Version B: reduced coordinate precision}

Export again. This time, in the \textbf{Layer Options} section of the
dialog, set \cd{COORDINATE\_PRECISION} to \cd{5}. Name it
\cd{districts\_B.geojson} and record the size.

\begin{tabularx}{\textwidth}{@{}c X@{}}
\toprule
\textbf{Decimal places} & \textbf{Resolves to} \\
\midrule
3 & 110 m \\
4 & 11 m \\
5 & 1.1 m \\
6 & 11 cm \\
\bottomrule
\end{tabularx}

\vspace{.8em}
Five decimal places locate a point to about a metre. A district boundary
digitised from 30 m imagery does not justify more than that, and writing
more both wastes bandwidth and overstates the accuracy of the data.

\subsection{Versions C and D: simplified geometry}

Use \textbf{Vector}, then \textbf{Geometry Tools}, then
\textbf{Simplify}. Run it twice on the original layer.

\begin{tabularx}{\textwidth}{@{}L{3cm} L{3.6cm} X@{}}
\toprule
\textbf{Version} & \textbf{Tolerance} & \textbf{Intended use} \\
\midrule
C & 0.001 degrees, about 100 m & A provincial or national view \\
D & 0.01 degrees, about 1 km   & A thumbnail or an overview map \\
\bottomrule
\end{tabularx}

\vspace{.8em}
Export each result to GeoJSON with precision 5, as
\cd{districts\_C.geojson} and \cd{districts\_D.geojson}. Record both
sizes.

\begin{keypoint}{Gaps along shared borders}
Simplifying adjacent polygons independently moves their shared borders
apart, leaving slivers and gaps. In the Simplify dialog, keep
\textbf{Preserve topology} enabled. If you see hairline gaps between
districts later, this is the cause.
\end{keypoint}


\subsection{Record your measurements}

Fill this in. These are your numbers, not the lecture's.

\vspace{.6em}
\begin{tabularx}{\textwidth}{@{}L{2.2cm} X r r@{}}
\toprule
\textbf{Version} & \textbf{What was done} & \textbf{Size} & \textbf{\% of A} \\
\midrule
Shapefile & All components together           & \hspace{2cm} & n/a \\
A         & Straight conversion               & & \\
B         & Precision 5                       & & \\
C         & Simplified 100 m, precision 5     & & \\
D         & Simplified 1 km, precision 5      & & \\
\bottomrule
\end{tabularx}

\vspace{1em}
Copy version C into your project as the one the map will use:

\begin{code}
copy districts_C.geojson C:\webgis-lab\data\districts.geojson
\end{code}

\begin{keypoint}{Checkpoint 2}
The completed measurement table, and a screenshot showing versions A and
D drawn together in QGIS so the loss of detail is visible.
\end{keypoint}


%==============================================================
\section{Read the data as text}

GeoJSON is a text format, which means you can inspect it directly. This
is the main practical advantage it has over a shapefile.

\subsection{Open it}

Open \cd{data/districts.geojson} in VS Code. Press
\cd{Shift+Alt+F} to format it, which puts each item on its own line and
makes the structure visible.

\begin{keypoint}{Format a copy, not the file you will serve}
Formatting inserts spaces and line breaks, which can add 30 percent to
the file size. Format it to read it, then undo with \cd{Ctrl+Z} before
saving.
\end{keypoint}


\subsection{Find the structure}

Answer these from the file itself:

\begin{enumerate}
  \item What is the \cd{type} of the outermost object?
  \item How many entries are in the \cd{features} array? VS Code shows
        the count if you select the array.
  \item Take the first feature. What is its geometry \cd{type}?
  \item How many levels of square bracket are there before the first
        number in its \cd{coordinates}?
  \item List the keys inside its \cd{properties}.
\end{enumerate}

\subsection{Check the coordinate order}

Find the first coordinate pair in the file.

\paragraph{What the numbers should look like}
Latitude in Pakistan runs from about 24 to 37. Longitude runs from about
61 to 77.

GeoJSON writes longitude first, so the first number of each pair should
be in the sixties or seventies, and the second in the twenties or
thirties. If they are the other way round, the file is wrong and every
feature will be drawn in the wrong place.


\subsection{Check that the rings close}

Take any polygon. Compare the first coordinate pair in its outer ring
with the last one. They must be identical. A ring that does not close is
invalid, and some renderers will reject the whole file rather than that
one feature.

\begin{keypoint}{Checkpoint 3}
Your five answers from 4.2, plus the first coordinate pair you found and
a sentence confirming the order is correct.
\end{keypoint}


%==============================================================
\section{Build the page}

\subsection{index.html}

\codecap{index.html, complete}
\begin{code}
<!DOCTYPE html>
<html lang="en">
<head>
  <meta charset="utf-8">
  <meta name="viewport" content="width=device-width, initial-scale=1">
  <title>District Explorer</title>

  <link rel="stylesheet"
        href="https://unpkg.com/leaflet@1.9.4/dist/leaflet.css">
  <link rel="stylesheet" href="style.css">
</head>
<body>

  <header>
    <h1>District Explorer</h1>
    <p class="subtitle">Districts of Pakistan, shaded by population</p>
  </header>

  <div id="panel">
    <label for="province">Province</label>
    <select id="province">
      <option value="all">All provinces</option>
    </select>
    <span id="count">loading...</span>
  </div>

  <div id="map"></div>

  <script src="https://unpkg.com/leaflet@1.9.4/dist/leaflet.js"></script>
  <script src="app.js"></script>
</body>
</html>
\end{code}

\paragraph{Why the script tags sit at the end of the body}
A script runs as soon as the browser reaches it. If \cd{app.js} ran
before the \cd{<div id="map">} existed, it could not find it, and
Leaflet would report that the map container was not found. Placing the
scripts last guarantees the elements exist first.

The Leaflet tag must come before \cd{app.js}, because your code uses
\cd{L}, which Leaflet defines.


\subsection{style.css}

\codecap{style.css, complete}
\begin{code}
body {
  font-family: system-ui, sans-serif;
  margin: 0;
  color: #23373b;
  background: #f4f6f6;
}

header {
  padding: 18px 24px 10px;
}

h1 {
  margin: 0;
  font-size: 24px;
}

.subtitle {
  margin: 4px 0 0;
  color: #5b6a6e;
}

#panel {
  padding: 10px 24px 14px;
  display: flex;
  align-items: center;
  gap: 12px;
}

#count {
  color: #5b6a6e;
  font-size: 14px;
}

#map {
  height: calc(100vh - 150px);
  width: 100%;
}

.legend {
  background: white;
  padding: 8px 10px;
  line-height: 20px;
  border-radius: 4px;
  box-shadow: 0 1px 4px rgba(0,0,0,.25);
  font-size: 13px;
}

.legend i {
  width: 16px;
  height: 16px;
  float: left;
  margin-right: 8px;
  opacity: .85;
}
\end{code}

\begin{keypoint}{The height rule}
\cd{\#map} must have a height. An empty \cd{<div>} has zero height, and
Leaflet will draw a map zero pixels tall into it with no error message
and a clean console.

\cd{calc(100vh - 150px)} means the full window height minus the space
taken by the header and panel. If your map does not appear, this line is
the first thing to check.
\end{keypoint}


\subsection{View it}

Right-click in the editor and choose \textbf{Open with Live Server}.
Confirm that the address bar reads \cdb{http://127.0.0.1:5500} and not
\cdb{file:///C:/...}.

\paragraph{Why this matters now rather than later}
In Part 6 the page will request a file with \cd{fetch}. A browser blocks
that request for pages opened directly from the file system. Opened that
way, your data will silently fail to load and the error will not explain
why.


\begin{keypoint}{Checkpoint 4}
A screenshot showing the header and the empty grey map area, with the
address bar visible on \cd{127.0.0.1}.
\end{keypoint}


%==============================================================
\section{Load the data with JavaScript}

\subsection{Create the map}

\codecap{app.js, first version}
\begin{code}
const map = L.map('map').setView([30.3753, 69.3451], 5);

L.tileLayer('https://tile.openstreetmap.org/{z}/{x}/{y}.png', {
  maxZoom: 19,
  attribution: '&copy; OpenStreetMap contributors'
}).addTo(map);
\end{code}

Save and reload. You should have a pannable map centred on Pakistan.

\paragraph{Coordinate order, again}
\cd{setView} takes \cd{[latitude, longitude]}. Your GeoJSON file writes
\cd{[longitude, latitude]}. The two conventions are both established and
neither is going to change, so the only defence is knowing which one you
are addressing.

Leaflet reads GeoJSON correctly on its own. It is only the coordinates
you type by hand that need care.


\subsection{Fetch the file}

Add this below the tile layer.

\codecap{app.js, continued}
\begin{code}
let districts = null;

async function loadData() {
  const response = await fetch('data/districts.geojson');

  if (!response.ok) {
    console.error('Request failed with status', response.status);
    return;
  }

  districts = await response.json();

  console.log('features:', districts.features.length);
  console.log('first feature:', districts.features[0]);
  console.log('properties:', districts.features[0].properties);
}

loadData();
\end{code}

Open the Console tab in the developer tools and reload.

\begin{enumerate}
  \item Confirm the feature count matches what QGIS reported.
  \item Expand \cd{first feature} and walk down to its coordinates.
  \item Read the property names from \cd{properties}, and confirm they
        match the ones you will use in the next part.
\end{enumerate}

\paragraph{What \texttt{await} is doing}
The file has to travel over the network, which takes time and can fail.
\cd{await} pauses until the response arrives. Without it, the next line
would run before the data existed, and \cd{districts} would still be
\cd{null}.

The \cd{response.ok} check catches the case where the server answered but
the file was not found. A 404 means the path is wrong; look at your own
code. A 500 means the server failed; look at the server.


\subsection{Try the array methods}

With the page loaded, type these straight into the Console. They operate
on real data and the results print immediately.

\begin{code}
districts.features.length

districts.features.map(f => f.properties.district).slice(0, 10)

districts.features.filter(f => f.properties.province === 'Punjab').length

[...new Set(districts.features.map(f => f.properties.province))].sort()
\end{code}

\begin{tabularx}{\textwidth}{@{}L{3.2cm} X@{}}
\toprule
\textbf{Method} & \textbf{Does} \\
\midrule
\cd{map}    & Turns every feature into something else, here its name \\
\cd{filter} & Keeps the features that pass a test \\
\cd{Set}    & Removes duplicates, giving the list of distinct provinces \\
\bottomrule
\end{tabularx}

\vspace{.8em}
The last line is the one you will use in Part 8 to fill the drop-down,
so note what it returns.

\begin{keypoint}{Checkpoint 5}
A screenshot of the Console showing the feature count, the expanded
first feature, and the list of distinct provinces.
\end{keypoint}


%==============================================================
\section{Draw it with Leaflet}

\subsection{Add the layer}

Add a \cd{draw} function and call it at the end of \cd{loadData}.

\codecap{app.js, continued}
\begin{code}
let layer = null;

function draw(features) {
  if (layer) {
    map.removeLayer(layer);
  }

  layer = L.geoJSON(
    { type: 'FeatureCollection', features: features },
    {
      style: {
        color: '#ffffff',
        weight: 1,
        fillColor: '#8c9ea3',
        fillOpacity: 0.75
      }
    }
  ).addTo(map);

  if (features.length > 0) {
    map.fitBounds(layer.getBounds());
  }

  document.getElementById('count').textContent =
    features.length + ' districts';
}
\end{code}

Then at the end of \cd{loadData}, after the \cd{console.log} lines:

\begin{code}
  draw(districts.features);
\end{code}

\noindent In Part 8 you will add one more line immediately above that
one, so the finished end of \cd{loadData} reads:

\begin{code}
  fillProvinceList();
  draw(districts.features);
\end{code}

\paragraph{Why the two guards are there}
\cd{removeLayer} runs first because \cd{L.geoJSON(...).addTo(map)} adds a
layer, it does not replace one. Without it, filtering would stack layers
on top of each other and the map would slow down with every change.

\cd{fitBounds} is skipped when the list is empty, because an empty layer
has no bounds and Leaflet throws \cd{Bounds are not valid}, which stops
the rest of the function from running.




\subsection{Check the network activity}

Open the \textbf{Network} tab, tick \textbf{Disable cache}, and reload.

\begin{enumerate}
  \item Find the request for \cd{districts.geojson}. Record its size and
        the time it took.
  \item Count the requests ending in \cd{.png}. These are map tiles.
  \item Click one tile and open its \textbf{Preview}. You are looking at
        a single 256 by 256 pixel square of map.
  \item Read the three numbers in that tile's URL and identify which is
        the zoom level.
  \item Pan the map and watch new tile requests arrive.
\end{enumerate}

\paragraph{Throttling the connection}
At the top of the Network tab there is a throttling drop-down, usually
set to \textbf{No throttling}. Set it to \textbf{Slow 3G} and reload.

The load time you now see is closer to what a user on a mobile
connection experiences. Record it. Then try the same with
\cd{districts\_A.geojson} substituted in, and record that too. Set the
throttle back to \textbf{No throttling} afterwards.


\begin{keypoint}{Checkpoint 6}
A screenshot of the Network tab showing the GeoJSON request with its
size, and a selected tile request with its Preview visible.
\end{keypoint}


%==============================================================
\section{Classify and add a popup}

\subsection{Shade by population}

Replace the fixed \cd{fillColor} with a function. Add this above
\cd{draw}.

\codecap{app.js, continued}
\begin{code}
function show(value) {
  if (value === null || value === undefined || value === '') {
    return 'not recorded';
  }
  return typeof value === 'number' ? value.toLocaleString() : value;
}

function colourFor(population) {
  if (population > 3000000) return '#7f2704';
  if (population > 1500000) return '#d94801';
  if (population >  750000) return '#f16913';
  if (population >  300000) return '#fd8d3c';
  return '#fdd0a2';
}

function styleFor(feature) {
  return {
    color: '#ffffff',
    weight: 1,
    fillColor: colourFor(feature.properties.pop_2023),
    fillOpacity: 0.8
  };
}
\end{code}

Then change the style option inside \cd{draw} to use it:

\begin{code}
      style: styleFor,
\end{code}

\paragraph{This is a classification}
The chain of \cd{if} statements is a manual classification with five
classes and four break values. Changing the breaks changes the map's
message, which is why the breaks belong in your report alongside the
map.

If every district comes out the same colour, inspect the values
themselves. A plain numeric string such as \cd{"1500000"} is converted to
a number by the comparison and classifies correctly. A value with a
thousands separator, a stray space, or a placeholder such as \cd{"NA"}
cannot be converted, so every comparison is false and the feature falls
through to the last line.


\subsection{Add a popup}

Add an \cd{onEachFeature} option to the \cd{L.geoJSON} call, directly
after \cd{style}.

\codecap{inside the options object in draw}
\begin{code}
      onEachFeature: function (feature, featureLayer) {
        const p = feature.properties;
        featureLayer.bindPopup(
          '<strong>' + show(p.district) + '</strong><br>' +
          'Province: ' + show(p.province) + '<br>' +
          'Population: ' + show(p.pop_2023) + '<br>' +
          'Area: ' + show(p.area_km2) + ' km<sup>2</sup>'
        );
      }
\end{code}

\paragraph{This is a callback}
You are not running that function. You are handing it to Leaflet, which
runs it once for every feature as the layer is built. The same idea
appears again in Part 8 with the drop-down.


\begin{keypoint}{Why the popup goes through \texttt{show}}
Calling \cd{.toLocaleString()} straight on an attribute fails if that
attribute is \cd{null}, or if the field is named differently in your copy
of the data. The failure is not limited to one feature: it happens while
Leaflet is building the layer, so the whole map fails to draw and the
console reports \cd{Cannot read properties of null}.

\cd{show} returns a readable placeholder instead, so one missing value
costs you one popup line rather than the entire map.
\end{keypoint}


\subsection{Add a legend}

\codecap{app.js, at the end}
\begin{code}
const legend = L.control({ position: 'bottomright' });

legend.onAdd = function () {
  const div = L.DomUtil.create('div', 'legend');
  const breaks = [0, 300000, 750000, 1500000, 3000000];
  const labels = ['under 300k', '300k to 750k', '750k to 1.5m',
                  '1.5m to 3m', 'over 3m'];

  for (let i = 0; i < breaks.length; i++) {
    div.innerHTML +=
      '<i style="background:' + colourFor(breaks[i] + 1) + '"></i> ' +
      labels[i] + '<br>';
  }
  return div;
};

legend.addTo(map);
\end{code}

\begin{keypoint}{Checkpoint 7}
A screenshot of the shaded map with the legend visible and one popup
open.
\end{keypoint}


%==============================================================
\section{Filter with a control}
\textit{Extension exercise}

\subsection{Fill the drop-down from the data}

Add this function, and call it from \cd{loadData} before \cd{draw}.

\codecap{app.js, continued}
\begin{code}
function fillProvinceList() {
  const select = document.getElementById('province');

  const provinces = [...new Set(
    districts.features.map(f => f.properties.province)
  )].sort();

  for (const name of provinces) {
    const option = document.createElement('option');
    option.value = name;
    option.textContent = name;
    select.appendChild(option);
  }
}
\end{code}

\paragraph{Building the list from the data}
The drop-down is not typed by hand. It is derived from the file, so it
stays correct if the data changes. This is the same \cd{map} and
\cd{Set} combination you tried in the Console in Part 6.


\subsection{React to the user}

\codecap{app.js, continued}
\begin{code}
document.getElementById('province')
  .addEventListener('change', function (event) {
    const chosen = event.target.value;

    const subset = chosen === 'all'
      ? districts.features
      : districts.features.filter(
          f => f.properties.province === chosen
        );

    draw(subset);
  });
\end{code}

Reload and change the drop-down. The map should redraw with only the
chosen province, zoom to it, and update the district count.

\begin{tabularx}{\textwidth}{@{}L{4.4cm} X@{}}
\toprule
\textbf{Piece} & \textbf{Role} \\
\midrule
\cd{getElementById}    & Finds the element in the page \\
\cd{addEventListener}  & Registers a function to run later \\
\cd{'change'}          & The event to wait for \\
\cd{event.target.value} & What the user chose \\
\cd{filter}            & Keeps only the matching features \\
\bottomrule
\end{tabularx}

\vspace{.8em}
Every interactive control in a dashboard is a variation on these five
lines: find an element, listen for an event, filter the data, redraw.

\begin{keypoint}{Checkpoint 8}
Two screenshots, one showing all provinces and one showing a single
province selected, with the district count visible in both.
\end{keypoint}


%==============================================================
\section{Measure the coordinate systems}
\textit{Extension exercise}

The lecture claimed that an area computed in Web Mercator at this
latitude is about 40 percent too large. This part tests that claim with
your own data.

\subsection{Produce two projected copies}

In QGIS, export the original layer twice.

\begin{tabularx}{\textwidth}{@{}L{4.2cm} L{3cm} X@{}}
\toprule
\textbf{Export as} & \textbf{CRS} & \textbf{Is} \\
\midrule
\cd{d\_3857}  & EPSG:3857  & Web Mercator, what the map is drawn in \\
\cd{d\_32643} & EPSG:32643 & UTM zone 43N, a projected system for this region \\
\bottomrule
\end{tabularx}

\subsection{Compute areas in each}

On each layer, open the \textbf{Field Calculator} and create a new
decimal field called \cd{area\_calc} with the expression:

\begin{code}
$area / 1000000
\end{code}

\cd{\$area} is computed in the layer's own coordinate system, so the two
layers will give different answers for the same district.

\subsection{Compare}

Pick five districts spread from the south of the country to the north.
Record their values.

\vspace{.6em}
\begin{tabularx}{\textwidth}{@{}X r r r r@{}}
\toprule
\textbf{District} & \textbf{Latitude} & \textbf{3857 km\textsuperscript{2}}
& \textbf{UTM km\textsuperscript{2}} & \textbf{\% difference} \\
\midrule
 & & & & \\
 & & & & \\
 & & & & \\
 & & & & \\
 & & & & \\
\bottomrule
\end{tabularx}

\subsection{Answer these}

\begin{enumerate}
  \item Does the percentage difference grow or shrink as you move north?
  \item The predicted inflation factor is $1/\cos^{2}(\mathrm{latitude})$.
        Compute it for one of your districts and compare it with the
        measured difference.
  \item Which of the two figures would you put in a report, and why?
\end{enumerate}

\begin{keypoint}{The working rule}
Store data in EPSG:4326. Display it in EPSG:3857, which the mapping
library handles for you. Measure in a projected system such as UTM zone
42N or 43N. Never report an area or a distance computed in Web Mercator.
\end{keypoint}


\begin{keypoint}{Checkpoint 9}
The completed comparison table and your three answers.
\end{keypoint}


%==============================================================
\section{Publish and submit}

\subsection{Commit and push}

\begin{code}
cd C:\webgis-lab
git add -A
git commit -m "District Explorer: map, classification, filter"
git push -u origin main
\end{code}

\paragraph{A browser window will open}
Git for Windows authenticates through your browser rather than asking
for a password in the terminal. Approve it once and it remembers. If
anything asks you to type your GitHub password into PowerShell, stop and
ask for help.


\subsection{Turn on hosting}

\begin{enumerate}
  \item On the repository page, open \textbf{Settings}, then
        \textbf{Pages}.
  \item Under Source, choose \textbf{Deploy from a branch}.
  \item Set the branch to \cd{main} and the folder to \cd{/ (root)}.
        Save.
  \item Wait one to three minutes, then refresh until the live URL
        appears.
\end{enumerate}

Open the URL on your phone, using mobile data rather than campus wifi.
The map is now being served over the public internet.

\subsection{Write the report}

Create \cd{report.md} in the repository, containing:

\begin{enumerate}
  \item Your nine checkpoint screenshots.
  \item The file size table from Part 3.
  \item The load times you measured with and without throttling.
  \item The classification breaks you used, and one sentence on why.
  \item The coordinate system comparison table from Part 9, if you
        completed it.
  \item Your answers to the questions below.
\end{enumerate}

\subsection{Questions}

\begin{enumerate}
  \item Version D is far smaller than version A. Give one situation in
        which version D would be the wrong choice.
  \item The page makes one request for your GeoJSON file and dozens for
        tiles. Explain the difference in how those two kinds of data
        reach the browser.
  \item Why does the map container need a height when the heading does
        not?
  \item \cd{draw} removes the existing layer before adding a new one.
        What would go wrong if it did not?
  \item A classmate's districts all appear in the same colour. List two
        possible causes and how you would tell them apart.
  \item You open the page by double-clicking \cd{index.html} instead of
        using Live Server. The basemap appears but the districts do not.
        Explain why.
\end{enumerate}

\subsection{Marking}

\begin{tabularx}{\textwidth}{@{}X r@{}}
\toprule
\textbf{Item} & \textbf{Marks} \\
\midrule
Published map that loads, with shading, legend and popups & 30 \\
File size measurement table, with correct method          & 15 \\
Network measurements, including the throttled figures     & 10 \\
Nine checkpoint screenshots                               & 15 \\
Answers to the six questions                              & 20 \\
Repository with a readable commit history                 & 10 \\
\midrule
Extensions in Parts 8 and 9                               & +10 bonus \\
\bottomrule
\end{tabularx}

\vspace{1em}
Submit two links on LMS: your repository and your live page.

%==============================================================
\section*{Appendix A: Troubleshooting}
\addcontentsline{toc}{section}{Appendix A: Troubleshooting}

\begin{tabularx}{\textwidth}{@{}L{4.6cm} X@{}}
\toprule
\textbf{Symptom} & \textbf{Cause and fix} \\
\midrule
Nothing where the map should be &
  The container has no height. Check the \cd{\#map} rule in
  \cd{style.css} \\
Grey map area, no tiles &
  No internet, or the tile host is blocked. Check that you can open
  \cd{unpkg.com} in a tab \\
\cd{L is not defined} &
  The Leaflet script tag is missing, misspelled, or placed after
  \cd{app.js}. Leaflet must load first \\
\cd{Map container not found} &
  The id in the HTML and the id in \cd{L.map()} do not match. Both must
  be exactly \cd{map} \\
Basemap loads, districts do not &
  Opened from the file system. Use Live Server. Check the address bar
  reads \cd{127.0.0.1} \\
404 on \cd{districts.geojson} &
  Wrong path. The file must be at \cd{data/districts.geojson} relative
  to \cd{index.html} \\
\cd{Cannot read properties of null} &
  \cd{draw} ran before the data arrived. Check that \cd{await} is
  present on both lines in \cd{loadData} \\
Every district the same colour &
  The values cannot be read as numbers: a thousands separator, a stray
  space, or a placeholder such as \cd{NA}. Check the attribute table \\
Map draws nothing and the console says \cd{Cannot read properties of
  null} &
  A popup is reading an attribute that is \cd{null} or named differently.
  Use the \cd{show} helper from Part 7 \\
\cd{Bounds are not valid} &
  \cd{fitBounds} was called on an empty selection. Check the
  \cd{features.length > 0} guard in \cd{draw} \\
Districts drawn in the ocean &
  Longitude and latitude are swapped in the source file \\
Hairline gaps between districts &
  Simplified without preserving topology. Re-run Simplify with that
  option enabled \\
Popup shows \cd{undefined} &
  The property name in the popup does not match the attribute table \\
Page works locally, not when published &
  A path with a capital letter or a backslash. GitHub Pages is
  case-sensitive and needs forward slashes \\
\bottomrule
\end{tabularx}

%==============================================================
\section*{Appendix B: Reference}
\addcontentsline{toc}{section}{Appendix B: Reference}

\subsection*{Commands}

\begin{code}
cd C:\webgis-lab          go to the project folder
code .                    open the folder in VS Code
git status                what has changed
git add -A                stage all changes
git commit -m "..."       save a snapshot
git push                  upload to GitHub
git log --oneline         see the commit history
\end{code}

\subsection*{Developer tools}

\begin{tabularx}{\textwidth}{@{}L{3.6cm} X@{}}
\toprule
\textbf{Tab} & \textbf{Use it for} \\
\midrule
Elements & Seeing the page as the browser built it, and editing CSS live \\
Console  & Reading errors, and printing values with \cd{console.log} \\
Network  & Request sizes, status codes, timings, and throttling \\
\bottomrule
\end{tabularx}

\subsection*{Numbers worth remembering}

\begin{tabularx}{\textwidth}{@{}L{6cm} X@{}}
\toprule
\textbf{Quantity} & \textbf{Value} \\
\midrule
Latitude range, Pakistan        & about 24 to 37 \\
Longitude range, Pakistan       & about 61 to 77 \\
Fifth decimal place             & about 1.1 m \\
Tile size                       & 256 by 256 pixels \\
EPSG:4326                       & WGS 84 degrees, for storage \\
EPSG:3857                       & Web Mercator metres, for display \\
EPSG:32642 and 32643            & UTM zones 42N and 43N, for measuring \\
Area inflation at 33.6\textdegree & about 1.4 times \\
\bottomrule
\end{tabularx}

\subsection*{Optional, if you finish everything}

\begin{enumerate}
  \item Add a second basemap and a layer control to switch between them.
  \item Highlight a district on mouseover by changing its outline weight.
  \item Add a text box that filters districts by name as the user types.
  \item Sort the drop-down by district count rather than alphabetically.
  \item Replace the manual classification breaks with quantiles computed
        from the data.
\end{enumerate}

\end{document}
